% ===== Math & Geometry (8 problems) =====
\section{Math \& Geometry {\normalsize\color{muted}$\cdot$ 8 problems}}
\markboth{Math \& Geometry}{}

\begin{problem}{48}{Rotate Image}{Medium}
\summary{Rotate an n x n matrix 90 degrees clockwise in place.}
\keyidea{Transpose, then reverse each row.}
\complexity{O(n\textasciicircum{}2) time, O(1) space}
\begin{code}
class Solution:
    def rotate(self, matrix: List[List[int]]) -> None:
        n = len(matrix)
        for i in range(n):
            for j in range(i + 1, n):
                matrix[i][j], matrix[j][i] = matrix[j][i], matrix[i][j]
        for row in matrix:
            row.reverse()
\end{code}
\end{problem}

\begin{problem}{54}{Spiral Matrix}{Medium}
\summary{Return all elements in clockwise spiral order.}
\keyidea{Shrink four boundaries (top, bottom, left, right) after each side is read.}
\complexity{O(m * n) time, O(1) extra space}
\begin{code}
class Solution:
    def spiralOrder(self, matrix: List[List[int]]) -> List[int]:
        res = []
        top, bottom, left, right = 0, len(matrix) - 1, 0, len(matrix[0]) - 1
        while top <= bottom and left <= right:
            for c in range(left, right + 1): res.append(matrix[top][c])
            top += 1
            for r in range(top, bottom + 1): res.append(matrix[r][right])
            right -= 1
            if top <= bottom:
                for c in range(right, left - 1, -1): res.append(matrix[bottom][c])
                bottom -= 1
            if left <= right:
                for r in range(bottom, top - 1, -1): res.append(matrix[r][left])
                left += 1
        return res
\end{code}
\end{problem}

\begin{problem}{73}{Set Matrix Zeroes}{Medium}
\summary{If a cell is 0, set its entire row and column to 0, in place.}
\keyidea{Use the first row/column as marker arrays; keep a separate flag for row 0.}
\complexity{O(m * n) time, O(1) space}
\begin{code}
class Solution:
    def setZeroes(self, matrix: List[List[int]]) -> None:
        R, C = len(matrix), len(matrix[0])
        row0_zero = any(matrix[0][c] == 0 for c in range(C))
        for r in range(1, R):
            for c in range(C):
                if matrix[r][c] == 0:
                    matrix[0][c] = 0
                    matrix[r][0] = 0
        for r in range(1, R):
            for c in range(1, C):
                if matrix[0][c] == 0 or matrix[r][0] == 0:
                    matrix[r][c] = 0
        if matrix[0][0] == 0:
            for r in range(R):
                matrix[r][0] = 0
        if row0_zero:
            for c in range(C):
                matrix[0][c] = 0
\end{code}
\end{problem}

\begin{problem}{202}{Happy Number}{Easy}
\summary{Repeatedly replace n with the sum of squares of its digits. Does it reach 1?}
\keyidea{Either reaches 1 or loops: detect the loop with a set (or Floyd).}
\complexity{O(log n) per step, O(log n) space}
\begin{code}
class Solution:
    def isHappy(self, n: int) -> bool:
        seen = set()
        while n != 1 and n not in seen:
            seen.add(n)
            n = sum(int(d) ** 2 for d in str(n))
        return n == 1
\end{code}
\end{problem}

\begin{problem}{66}{Plus One}{Easy}
\summary{Add one to a number stored as a list of digits.}
\keyidea{From the end: 9 becomes 0 and carry; otherwise increment and return. All 9s -\textgreater{} prepend 1.}
\complexity{O(n) time, O(1) space}
\begin{code}
class Solution:
    def plusOne(self, digits: List[int]) -> List[int]:
        for i in range(len(digits) - 1, -1, -1):
            if digits[i] < 9:
                digits[i] += 1
                return digits
            digits[i] = 0
        return [1] + digits
\end{code}
\end{problem}

\begin{problem}{50}{Pow(x, n)}{Medium}
\summary{Compute x raised to integer n (n may be negative).}
\keyidea{Fast exponentiation: square x and halve n; multiply in x when the low bit is 1.}
\complexity{O(log n) time, O(1) space}
\begin{code}
class Solution:
    def myPow(self, x: float, n: int) -> float:
        if n < 0:
            x, n = 1 / x, -n
        res = 1.0
        while n:
            if n & 1:
                res *= x
            x *= x
            n >>= 1
        return res
\end{code}
\end{problem}

\begin{problem}{43}{Multiply Strings}{Medium}
\summary{Multiply two non-negative numbers given as strings, without converting to int.}
\keyidea{Grade-school multiplication: digit i times digit j lands at position i + j (on reversed strings).}
\complexity{O(m * n) time, O(m + n) space}
\begin{code}
class Solution:
    def multiply(self, num1: str, num2: str) -> str:
        if num1 == "0" or num2 == "0":
            return "0"
        res = [0] * (len(num1) + len(num2))
        a, b = num1[::-1], num2[::-1]
        for i in range(len(a)):
            for j in range(len(b)):
                res[i + j] += int(a[i]) * int(b[j])
                res[i + j + 1] += res[i + j] // 10
                res[i + j] %= 10
        while len(res) > 1 and res[-1] == 0:
            res.pop()
        return "".join(map(str, res[::-1]))
\end{code}
\end{problem}

\begin{problem}{2013}{Detect Squares}{Medium}
\summary{add(point) to a multiset; count(point) returns how many axis-aligned squares it completes with stored points.}
\keyidea{Treat each stored point as the diagonal corner; if it forms a square, multiply counts of the other two corners.}
\complexity{add O(1), count O(n) time; O(n) space}
\begin{code}
class DetectSquares:
    def __init__(self):
        self.pts = Counter()

    def add(self, point: List[int]) -> None:
        self.pts[tuple(point)] += 1

    def count(self, point: List[int]) -> int:
        px, py = point
        res = 0
        for (x, y), c in list(self.pts.items()):
            if abs(px - x) != abs(py - y) or x == px or y == py:
                continue
            res += c * self.pts[(x, py)] * self.pts[(px, y)]
        return res
\end{code}
\end{problem}
